\documentclass[11pt,a4paper]{article}
\usepackage[T1]{fontenc}
\usepackage{lmodern,amsmath,amssymb,amsthm,mathtools,booktabs,microtype}
\usepackage[margin=27mm]{geometry}
\usepackage[colorlinks=true,linkcolor=blue,citecolor=blue,urlcolor=blue]{hyperref}
\setlength{\emergencystretch}{2em}
\newtheorem{theorem}{Theorem}[section]
\newtheorem{lemma}[theorem]{Lemma}
\newtheorem{proposition}[theorem]{Proposition}
\newtheorem{corollary}[theorem]{Corollary}
\newcommand{\Q}{\mathbb Q}\newcommand{\Z}{\mathbb Z}
\DeclareMathOperator{\Tr}{Tr}\DeclareMathOperator{\Gal}{Gal}
\title{An obstruction to finite M\"obius coupling of six elliptic contacts}
\author{}\date{17 November 2026}
\begin{document}\maketitle
\begin{abstract}
The six-contact polynomial of the conditional elliptic source has
exactly twelve rational M\"obius symmetries. Coupling the two source
arguments by any of these symmetries does not preserve the prescribed-index
completeness theorem. For input one, ten branches have no rational slope
for either retained power scale and every permitted coefficient shift.
The identity violates a contact guard. The remaining involution has
simple contacts at inert primes for an unbounded progression of shifts,
forcing a negative odd detector pole for every auxiliary tuple.
Chebotarev gives such primes outside every fixed finite exceptional set.
The symmetry and pole-class choices may vary with the input and shift,
provided the pole classes retain the stated common local implication.
\end{abstract}

\section{The prescribed source and the statement}

Set
\[
 K=\Q(\sqrt2),\quad D=(3-\sqrt2)^2,\quad E:y^2=x^3-Dx,
 \quad P=\left(\frac{57-28\sqrt2}{8},
                    \frac{-660+417\sqrt2}{32}\right).
\]
The conditional source uses points $A,H,U$, the functions
$f=1/x$, $z=-x/y$, and the six field-valued tests
\begin{equation}\label{eq:tests}
\begin{gathered}
 r=\frac{z(U)}{z(H)},\qquad s=\frac{z(U-A)}{z(H-P)},\\
 C_1=f(H),\quad C_2=f(U),\quad
 C_3=\overline{f(H-P)},\quad C_4=\overline{f(U-A)},\\
 C_5=r-\bar s,\qquad C_6=r+\bar s-2X.
\end{gathered}
\end{equation}
All point, difference and quotient guards are retained. For a fixed odd
integer $B\ge3$ and $e\in\{4,6\}$, the two source arguments are
\begin{equation}\label{eq:source}
 T=\Tr x(A),\qquad a=(X^2+B)^e+T^e,
 \qquad t_0=\lambda,\quad t_1=a\lambda,
\end{equation}
and the contact polynomial is
\begin{equation}\label{eq:F}
 F(t)=\frac{t(t+1)(2t+1)(3t+1)(4t+1)(5t+2)}{120}.
\end{equation}
The literal source in~\cite{Reweighted} has
\begin{equation}\label{eq:detector}
 d_j=F(t_j),\quad w_j=\frac{d_j}{4(t_j-1)^8},\quad
 W_j=\prod_{i=1}^6\left(1+\frac{C_i^2}{d_j}\right),\quad
 V_j=\Tr(W_j^{-1}),\quad b_j=\frac{4w_j}{V_j^2}.
\end{equation}
The guards include $d_j(t_j-1)V_j\ne0$ and invertibility of $W_j$.
An intrinsic pole block for $7b_j$, for example, implies
\begin{equation}\label{eq:localpole}
 v_p(b_j)\ \text{is not a negative odd integer}\qquad(p\notin S),
\end{equation}
where $S$ is fixed and finite~\cite{Intrinsic}.

Write
\[
 \mathcal R=\{0,-1,-1/2,-1/3,-1/4,-2/5\},\qquad
 \mathcal M=\{\phi\in\operatorname{PGL}_2(\Q):\phi(\mathcal R)=\mathcal R\}.
\]
The coupling considered here is
\begin{equation}\label{eq:coupling}
                         a\lambda=\phi(\lambda),\qquad\phi\in\mathcal M.
\end{equation}
Every chosen M\"obius denominator is required to be nonzero.

\begin{theorem}\label{thm:main}
Fix $e\in\{4,6\}$ and a finite set $S$ containing $2,3,5,7$.
At the prescribed integer-input point
\[
                        X=1,\qquad A=P,
\]
there are infinitely many odd shifts $B\ge3$ for which no guarded
rational auxiliary tuple satisfies~\eqref{eq:source}--\eqref{eq:coupling}
and~\eqref{eq:localpole}, even if only one of the two pole conditions
is retained. The index of that pole condition and the symmetry
$\phi$ may be chosen separately for each tuple and shift.

The assertion also holds if $d_j$ in~\eqref{eq:detector} is replaced
by $\varepsilon_jF(t_j)$, with $\varepsilon_j$ chosen from a fixed finite
subset of $\Q^\times$. It suffices to enlarge $S$ by the numerator
and denominator primes of these constants. Pole predicates may be
chosen from any family whose members all imply~\eqref{eq:localpole}
outside the same fixed $S$.
\end{theorem}

Thus no choice from the complete symmetry menu supplies the prescribed
completeness assertion for every sufficiently large odd shift. The
theorem concerns the displayed source, height scales and prescribed
root point. It does not exclude a particular fixed shift, a changed
height scale, an independently chosen root index, or another contact
system. It does not cover variable nonunit rescalings that change the
local detector condition or derivative detectors with a different
contact polynomial. If a restricted source retains the source's
soundness conclusion that the root index must equal $X$, the theorem
also obstructs that source's eventual completeness.

\section{The complete symmetry menu}

\begin{lemma}\label{lem:menu}
The group $\mathcal M$ has precisely the following twelve elements.
For $\phi(t)=(\alpha t+\beta)/(\gamma t+\delta)$, the last column is
the constant in
\[
                F(\phi(t))(\gamma t+\delta)^6=\kappa F(t).
\]
\begin{center}
\begin{tabular}{rrrrr}\toprule
$\alpha$&$\beta$&$\gamma$&$\delta$&$\kappa$\\\midrule
1&0&0&1&1\\
1&0&$-5$&$-1$&$-1$\\
1&1&2&$-1$&27\\
4&1&$-7$&$-1$&$-27$\\
2&1&$-3$&$-2$&$-1$\\
3&1&$-7$&$-2$&1\\
2&1&$-7$&$-3$&1\\
3&1&$-8$&$-3$&$-1$\\
1&1&$-7$&$-4$&$-27$\\
4&1&$-13$&$-4$&27\\
5&2&$-11$&$-5$&27\\
5&2&$-14$&$-5$&$-27$\\\bottomrule
\end{tabular}
\end{center}
\end{lemma}
\begin{proof}
A M\"obius map is determined by the images of the three distinct
points $0,-1,-1/2$. There are $6\cdot5\cdot4=120$ possible ordered
images in $\mathcal R$. For each choice $(u_1,u_2,u_3)$, solve
\[
 \alpha r_i+\beta-u_i\gamma r_i-u_i\delta=0
 \qquad(r_1,r_2,r_3)=(0,-1,-1/2).
\]
The nullspace is one-dimensional. Discard singular matrices, maps
having a pole at a root, and maps whose six images are not exactly
$\mathcal R$. The retained primitive integer matrices are exactly
the twelve displayed ones. These finite rational computations and
each multiplier identity are reproduced by the exact checker.
\end{proof}

\section{Uniform square gaps for ten branches}

At $X=1,A=P$, the actual source trace and scale are
\begin{equation}\label{eq:actualscale}
 T=\frac{57}{4},\qquad
 a=\frac{p_B}{q},\qquad
 q=4^e,\quad p_B=q(B+1)^e+57^e.
\end{equation}
In particular $a>4$, and $p_B$ is odd.

The identity symmetry forces $(a-1)\lambda=0$, contradicting
$F(\lambda)\ne0$. For every row except the first two, coupling gives
\[
 a\gamma\lambda^2+(a\delta-\alpha)\lambda-\beta=0.
\]
A rational solution would require
\[
 G=(\delta p_B-\alpha q)^2+4\gamma\beta p_Bq
\]
to be an integer square. Set
\begin{equation}\label{eq:squaregap}
 N=\delta^2p_B+(-\delta\alpha+2\gamma\beta)q,
 \qquad K_0=4\beta\gamma(\alpha\delta-\beta\gamma).
\end{equation}
Direct expansion gives
\[
                           \delta^2G=N^2+K_0q^2.
\]

\begin{lemma}\label{lem:gap}
For all ten rows under consideration, both exponents, and every
integer $B\ge3$, one has
\[
                         N>0,\qquad 0<|K_0|q^2<2N-1.
\]
Consequently none of these ten branches has a rational slope.
\end{lemma}
\begin{proof}
The integer $N$ is increasing with $B$. At $B=3$, the respective
values of $p_B$ are $10621537$ and $34313224465$.
The following table gives $K_0$ and the positive gap
$2N-1-|K_0|q^2$ at that minimum; the rows follow Lemma~\ref{lem:menu}.
\begin{center}
\begin{tabular}{rrrr}\toprule
Row&$K_0$&Gap, $e=4$&Gap, $e=6$\\\midrule
3&$-24$&19672769&68223836705\\
4&$-84$&15732929&67217080865\\
5&12&84184839&274304452743\\
6&$-28$&83133191&274035968135\\
7&$-28$&189348561&617168212785\\
8&32&189086929&617101112113\\
9&$-84$&334379039&1096613814815\\
10&156&329660447&1095405855263\\
11&264&513765617&1711231882577\\
12&$-336$&509040881&1710023824721\\\bottomrule
\end{tabular}
\end{center}
Each $N$ is positive at the minimum. If $K_0<0$, the value
$N^2+K_0q^2$ lies strictly between $(N-1)^2$ and $N^2$.
If $K_0>0$, it lies strictly between $N^2$ and $(N+1)^2$.
It is therefore not an integer square, whereas a square $G$ would
make $\delta^2G$ one.
\end{proof}

\section{The involution and the inert-place obstruction}

The only remaining symmetry is $\phi(t)=-t/(5t+1)$. Its guarded
coupling has the unique solution
\begin{equation}\label{eq:involution}
 \lambda=-\frac{a+1}{5a},\qquad
 H(a)=(a^2-1)(4a-1)(3a-2)(2a-3)(a-4),
\end{equation}
with identities
\begin{equation}\label{eq:H}
                  -F(\lambda)=\frac{H(a)}{375000a^6},
                  \qquad F(a\lambda)=\frac{H(a)}{375000}.
\end{equation}

\begin{lemma}\label{lem:inert}
Let $p\ne2,7$ be inert in $K$. Suppose $v_p(d)=v_p(w)=1$ in
a defined detector
\[
 W=\prod_i(1+C_i^2/d),\qquad V=\Tr(W^{-1}),\qquad b=4w/V^2,
\]
whose tests include $f(H)$ and $\overline{f(H-P)}$.
Then $V=0$ or $v_p(b)$ is negative and odd.
\end{lemma}
\begin{proof}
The curve has good reduction at $p$: its discriminant is $64D^3$
and $\operatorname{N}_{K/\Q}(D)=49$. The displayed coordinates of
$P$ are integral, so its reduction is an affine point and is not
the identity. Conjugation fixes the unique place above an inert prime.
If both indicated tests had positive valuation, both $H$ and $H-P$
would reduce to the identity, forcing $P$ to do so. Thus one of these
tests has nonpositive valuation.

Every factor has valuation
\[
                         \min(0,2v(C_i)-1)\le0;
\]
the two summands have different valuation parity. A zero test gives
factor one. At least one factor has valuation at most $-1$.
Hence $W^{-1}\in p\mathcal O_{K_p}$ and $V\in p\Z_p$. If $V\ne0$,
then $v_p(b)=1-2v_p(V)$ is negative and odd.
\end{proof}

\section{Explicit congruence certificates}

The prime $p=76253$ satisfies $p>76176$ and $p\equiv5\pmod8$,
so it is inert in $K$. The following two shifts satisfy
$v_p(H(a))=1$ for the exact scale~\eqref{eq:actualscale}:
\begin{center}
\begin{tabular}{rrrl}\toprule
$e$&$B$&$a\pmod p$&Exact $a$\\\midrule
4&88285&4&$15552754887321993945697/256$\\
6&99329&$2/3$&$3934074326572091195237078520447249/4096$\\\bottomrule
\end{tabular}
\end{center}
The same valuation holds on the full odd progressions
\begin{equation}\label{eq:progressions}
 \begin{array}{ll}
 B\equiv88285\pmod{11629040018},&e=4,\\
 B\equiv99329\pmod{11629040018},&e=6,
 \end{array}
\end{equation}
because $11629040018=2p^2$ and the denominator of the polynomial
$a(B)$ is a unit at $p$. Each displayed contact root of $H$ is simple;
$a,\lambda,\lambda-1,a\lambda-1,5\lambda+1$ are units at $p$.
Equations~\eqref{eq:H} give valuation one for both contact values and
both physical weights. Lemma~\ref{lem:inert} therefore applies to
either detector, and to either sign of its contact function.

The finite certificate checks primality by division by every possible
divisor at most $276$, verifies Euler's criterion for $2$, and evaluates
the scale and contact valuations using exact rational arithmetic.
As a supplementary order certificate, write
$\mathbb F_{p^2}=\mathbb F_p[s]/(s^2-2)$. Then
\[
 \widetilde P=((28602,38123),(9511,26225))
\]
in the basis $1,s$, and
\[
 \operatorname{ord}(\widetilde P)=16518755=5\cdot11\cdot157\cdot1913.
\]
The checker verifies annihilation by this integer and nonannihilation
by its quotient by each of the four prime factors, entirely in the
finite field. The proof of Lemma~\ref{lem:inert} requires only the
nonidentity reduction, not its exact order.

\section{Primes outside every finite exceptional set}

\begin{lemma}\label{lem:chebotarev}
For each $e\in\{4,6\}$, infinitely many rational primes inert in
$\Q(\sqrt2)$ admit a simple residue root of
\[
                         Y^e+57^e-4^{e+1}.
\]
\end{lemma}
\begin{proof}
Put $m_e=57^e-4^{e+1}$. The exact factorizations are
\[
 m_4=10554977=17\cdot193\cdot3217,\qquad
 m_6=34296430865=5\cdot47\cdot3943\cdot37013.
\]
Thus $Y^4+m_4$ is Eisenstein at $17$, and $Y^6+m_6$ is
Eisenstein at $5$. Let $L_e=\Q(\alpha_e)$ for a root $\alpha_e$.
In each case the Eisenstein prime has a unique prime above it in
$L_e$, with total ramification and residue degree one.

The field $L_4$ cannot contain $\Q(\sqrt2)$: $17$ splits in that
quadratic field, as $6^2\equiv2\pmod{17}$, and the two primes
would have distinct primes above them in $L_4$. The field $L_6$
cannot contain $\Q(\sqrt2)$ either: $5$ is inert in that field,
so the residue degree of a prime of $L_6$ above $5$ would be
divisible by two.

Let $M_e$ be a Galois closure of $L_e(\sqrt2)$ over $\Q$.
Since $\sqrt2\notin L_e$, there is an automorphism
$\sigma\in\Gal(M_e/L_e)$ that sends $\sqrt2$ to $-\sqrt2$.
Chebotarev's theorem~\cite[Theorem 8.31]{Milne} supplies infinitely
many unramified rational primes whose Frobenius conjugacy class is
that of $\sigma$. Its action on $\sqrt2$ makes each prime inert
in the quadratic field. Since $\sigma$ fixes $\alpha_e$, its
permutation of the roots has a fixed point; after omitting the finitely
many polynomial-discriminant primes, the polynomial has a simple
root modulo each such prime.
\end{proof}

\begin{proof}[Proof of Theorem~\ref{thm:main}]
Enlarge $S$ to contain the numerator and denominator primes of the
allowed fixed multipliers. Choose a prime from
Lemma~\ref{lem:chebotarev} outside $S$ and the remaining finite
exceptional set, and a simple root $y$ modulo $p$.
The residue $B+1=y/4$ gives
\[
                      (B+1)^e+(57/4)^e\equiv4\pmod p.
\]
The derivative $e(B+1)^{e-1}$ is nonzero after omitting primes
dividing $em_e$. Also $H'(4)=11250$, so $H(a(B))$ has a simple
zero in this residue class. Of its $p$ lifts modulo $p^2$, precisely
one is a zero modulo $p^2$. Any other lift, chosen odd, gives
\[
                             v_p(H(a(B)))=1
\]
on an infinite progression modulo $2p^2$.

For $B\ge3$, Lemma~\ref{lem:gap} excludes ten possible symmetries
and the identity violates the contact guard. The involution is
therefore the only remaining candidate. At its residue $a=4$,
all the denominator guards in~\eqref{eq:H} and the physical weights
are units. Multiplication by any allowed $\varepsilon_j$ preserves
valuation one. Lemma~\ref{lem:inert} contradicts either retained
pole condition. Since this argument excludes all twelve possibilities
at every shift in the progression, allowing their choices to vary
does not change the conclusion.
\end{proof}

\section{Reproduction}

The exact computation is distributed at
\nolinkurl{negative_results/computation/six_contact_mobius_v62/}.
Run \texttt{python check.py --output certificate.json} in that directory.
The checker enumerates all $120$ possible ordered image triples,
verifies the twelve rational identities, the twenty uniform square-gap
inequalities, both explicit progressions, the finite-field order,
and the two Eisenstein factorizations. The independent standard-library
checker \texttt{check\_independent.py} reconstructs the symmetry group
by explicit interpolation, verifies identities by polynomial coefficient
multiplication, and checks the square gaps and the integral contact
numerators. The infinite assertion uses the
proof of Lemma~\ref{lem:chebotarev}; it is not inferred from a finite search.

\begin{thebibliography}{9}
\bibitem{Reweighted}
\emph{Quadratic reweightings of native elliptic detectors}, repository source,
\nolinkurl{papers/pole_parameters/reweighted_elliptic_detectors_20261116_v53.tex}.
\bibitem{Intrinsic}
\emph{A two-coordinate pole test of coefficient height four}, repository source,
\nolinkurl{papers/pole_parameters/intrinsic_pole_20261117_v56.tex}.
\bibitem{Milne}
J.~S.~Milne, \emph{Algebraic Number Theory}, course notes, Theorem~8.31,
\url{https://www.jmilne.org/math/CourseNotes/ANT.pdf}, accessed 17 November 2026.
\end{thebibliography}
\end{document}
